% Procedencia íntegra: HIBRIDACION_ACCION_MASA_PROPAGACION.md, §§2--7.
% La gestión, el historial y los localizadores se conservan en gestion/.
\section{Acción, masa y propagación: compatibilidad estructural y transporte}
\label{md:hib:seccion}

\subsection{Sección de acción compartida y ley de masas}
\label{md:hib:accion-masas}

Sean
\begin{equation}
 s_0=10^{-34}\mathcal U_S,\qquad
 \hbar_{\mathrm{ret}}=\eta_{\mathrm{ret}}s_0,\qquad
 R_{\mathrm{act}}=\frac{H_5}{\eta_{\mathrm{ret}}}.
 \label{md:hib:accion}
\end{equation}
Se mantienen las coordenadas construidas en grados:
\begin{equation}
 A_{\deg}=1000\alpha,\qquad
 C^*_{\deg}=2(\eta_{\mathrm{ret}}+\alpha).
 \label{md:hib:angulos}
\end{equation}
Sólo después se transportan a radianes mediante
$x=\pi A_{\deg}/180$ e $y=\pi C^*_{\deg}/180$. Los canales son
$q_\pm=\exp(-x\mp y)$, con $x>|y|$. La forma elíptica se expresa por
\begin{equation}
 \zeta=\operatorname{artanh}(y/x).
 \label{md:hib:forma}
\end{equation}
La composición utiliza primero el coeficiente adimensional
$\eta_{\mathrm{ret}}$ y conserva así los tipos de acción y de ángulo.
En esta notación, $\eta_{\mathrm{ret}}$ es el coeficiente de acción,
$\eta_h(\Gamma)$ pertenece a la torre de una ruta y $\zeta$ expresa la
deformación de la elipse; sus funciones se mantienen separadas.

Se trabaja con $t_0,c_{\mathrm{int}},\eta_{\mathrm{ret}}>0$, carácter
convergente y factor $\Xi_\Gamma>0$. El lector másico construido determina
\begin{equation}
 \omega_0=\frac{2\pi}{108t_0},\qquad
 E_{\mathrm{clk}}=\hbar_{\mathrm{ret}}\omega_0,\qquad
 m_{\mathrm{clk}}=\frac{E_{\mathrm{clk}}}{c_{\mathrm{int}}^2}.
 \label{md:hib:reloj}
\end{equation}
Para una ruta con realización escalar,
\begin{equation}
 m_\Gamma^\beta=m_{\mathrm{clk}}\Xi_\Gamma,\qquad
 \Xi_\Gamma=b_e\,
 \chi_{\mathrm{ref}}(\sigma_\Gamma-\sigma_e,\eta_\Gamma-\eta_e)
 R_{\mathrm{act}}^{\kappa_\Gamma}.
 \label{md:hib:factor-ruta}
\end{equation}
$\Xi_\Gamma$ abrevia el factor estructural de ruta. Se conserva expresamente
\begin{equation}
 b_e=\frac{\varepsilon_e^{(0)}\mathcal U_E}{E_{\mathrm{clk}}}.
 \label{md:hib:normalizacion-electron}
\end{equation}
El estado central conserva su normalización propia respecto del cuanto
cronológico:
\begin{equation}
 \varepsilon_e^{(0)}=\frac{\sqrt3}{4}
 \left[\exp\!\left(\frac{\varphi}{\pi^2}\right)-22\alpha^3\right]
 (1+15\Delta_4),\qquad
 \varepsilon_e^\beta=\varepsilon_e^{(0)}
 R_{\mathrm{act}}^{80-54\alpha+6\Delta_4}.
 \label{md:hib:electron}
\end{equation}
El carácter refinado incluye los términos
\begin{equation}
 n_Ax+\frac{s_C}{6}y+k_{\Delta_4}\Delta_4+b_\zeta\zeta_{270}
 +\sum_h N_h(q_-^h-q_+^h),\qquad
 N_h=\eta_h+b_{120}\mathbf 1_{\{h=120\}},
 \label{md:hib:caracter}
\end{equation}
con $h$ en el semigrupo generado por $90$ y $120$. El término de altura
$120$ se incluye una vez. Los coeficientes primitivos de retorno se
construyen mediante
\begin{equation}
 \begin{aligned}
 a_n(\Gamma)&=\operatorname{Tr}(R_\Gamma U_\Gamma^n),\\
 p_n(\Gamma)&=\frac1n\sum_{d\mid n}\mu_{\mathrm{Mob}}(d)
                         a_{n/d}(\Gamma),\\
 \eta_h(\Gamma)&=p_h(\Gamma)-p_h(\Gamma_e),
 \end{aligned}
 \label{md:hib:retornos-primitivos}
\end{equation}
donde $\mu_{\mathrm{Mob}}$ es la función de Möbius aritmética.
La acción participa en la escala absoluta, en la razón $R_{\mathrm{act}}$
y en $y$, que determina los canales. En consecuencia, la cancelación de
una escala dimensional conserva las restantes dependencias estructurales.
Al variar el reloj, $b_e=\varepsilon_e^{(0)}\mathcal U_E/E_{\mathrm{clk}}$
varía con él: conservar la lectura electrónica exige transportar el
carácter completo, en vez de fijar por separado ese cociente.
En la realización longitudinal integrada se utiliza
$t_0=\pi_{\rm HMT}L_\alpha/(54c_{\mathrm{int}})$ de
\eqref{md:rigidez:reloj}.

\subsubsection{Frecuencia y longitud de cada sección masiva}
\label{md:hib:frecuencia-longitud}

Definiendo $E_\Gamma=m_\Gamma c_{\mathrm{int}}^2$ como energía de reposo,
$\Omega_\Gamma=E_\Gamma/\hbar_{\mathrm{ret}}$,
$\nu_\Gamma=\Omega_\Gamma/(2\pi)$ y
$R_{108}=c_{\mathrm{int}}/\omega_0$, se obtiene
\begin{equation}
 \begin{aligned}
 E_\Gamma&=\hbar_{\mathrm{ret}}\omega_0\Xi_\Gamma,&
 \Omega_\Gamma&=\omega_0\Xi_\Gamma,\\
 \nu_\Gamma&=\frac{\Xi_\Gamma}{108t_0},&
 \bar\lambda_\Gamma&=
 \frac{\hbar_{\mathrm{ret}}}{m_\Gamma c_{\mathrm{int}}}
 =\frac{R_{108}}{\Xi_\Gamma}.
 \end{aligned}
 \label{md:hib:frecuencias}
\end{equation}
\begin{proof}
Se sustituye $m_\Gamma=m_{\mathrm{clk}}\Xi_\Gamma$ y se efectúan las
cancelaciones en una misma coordenatización dimensional. Para rutas positivas $X,Y$,
la propiedad de carácter y la normalización común dan
\begin{equation}
 \frac{\Omega_Y}{\Omega_X}=\frac{m_Y}{m_X}
 =\chi_{\mathrm{ref}}(\sigma_Y-\sigma_X,\eta_Y-\eta_X)
 R_{\mathrm{act}}^{\kappa_Y-\kappa_X}.
 \label{md:hib:cocientes}
\end{equation}
\end{proof}
La cancelación es dimensional; permanecen los canales y la memoria
relativa. El sector nulo queda fuera del dominio del inverso de masa.
En una fibra no escalar se conserva el operador o su medida espectral,
en lugar de escoger un representante puntual arbitrario.

\subsubsection{Frecuencia levantada y compatibilidad de normalizaciones}
\label{md:hib:tres-frecuencias}

La construcción cronológica también produce la frecuencia levantada de
holonomía
\begin{equation}
 \omega_j=\frac{\theta_j+2\pi w_j}{108t_0},\qquad
 \lambda_j=r_j\exp(i\theta_j),
 \label{md:hib:frecuencia-levantada}
\end{equation}
donde $w_j$ conserva el devanado de la historia. Su lectura energética es
\begin{equation}
 E_j=\hbar_{\mathrm{ret}}\omega_j
 \chi_{\mathrm{full}}R_{\mathrm{act}}^{\kappa_j}.
 \label{md:hib:energia-levantada}
\end{equation}
Esta frecuencia, $\omega_0$ y $\Omega_\Gamma$, corresponden a lectores
distintos. Si ambas expresiones energéticas representan el mismo modo y
el mismo exponente, su compatibilidad exacta exige
\begin{equation}
 \frac{\omega_j}{\omega_0}\chi_{\mathrm{full}}
 =b_e\chi_{\mathrm{ref,rel}}.
 \label{md:hib:compatibilidad-normalizaciones}
\end{equation}
Las dos normalizaciones se identifican mediante esta igualdad y no se
multiplican entre sí. La fase reducida $\theta_j$ por sí sola no recupera
$w_j$. Por ello, un retorno de fase conserva una distinción con la historia
cronológica completa. El devanado conserva su constructor; esta distinción
no autoriza elegirlo libremente ni producir energía sin una dinámica.

\subsection{Reciprocidad espacial y composición gravitatoria}
\label{md:hib:gravedad}

La sección gravitatoria circular expresa el resultado del
teorema~\ref{md:rigidez:teorema}:
$G_a=c_{\mathrm{int}}^3L_\alpha^2/\hbar_a$ y $R_{108}=L_\alpha$.
La longitud procede del retorno inscrito y precede a esta
identificación circular. Manteniendo
$\hbar_a=\hbar_{\mathrm{ret}}$ y la convención reducida
$r_g=Gm/c_{\mathrm{int}}^2$, la masa anterior determina
\begin{equation}
 \boxed{r_{g,\Gamma}=R_{108}\Xi_\Gamma,\qquad
 \bar\lambda_\Gamma=\frac{R_{108}}{\Xi_\Gamma},\qquad
 r_{g,\Gamma}\bar\lambda_\Gamma=R_{108}^2.}
 \label{md:hib:reciprocidad-radial}
\end{equation}
\begin{proof}
La sustitución de $G_a$ y de $m_\Gamma$ da
\[
 \frac{G_am_\Gamma}{c_{\mathrm{int}}^2}
 =\frac{c_{\mathrm{int}}^3R_{108}^2}{\hbar_{\mathrm{ret}}}
 \frac{\hbar_{\mathrm{ret}}\omega_0\Xi_\Gamma}
 {c_{\mathrm{int}}^4}
 =R_{108}\Xi_\Gamma.
\]
La segunda igualdad procede de \eqref{md:hib:frecuencias}; su producto
prueba la tercera.
\end{proof}
La misma coordenada de ruta aumenta una lectura radial y contrae la otra.
En logaritmos,
\begin{equation}
 \log\frac{r_{g,\Gamma}}{R_{108}}=\log\Xi_\Gamma,
 \qquad
 \log\frac{\bar\lambda_\Gamma}{R_{108}}=-\log\Xi_\Gamma.
 \label{md:hib:radios-logaritmicos}
\end{equation}
El centro geométrico de ambas lecturas queda fijado por el área $R_{108}^2$.

En la versión operatoria acotada, sea $X$ autoadjunto y acotado, con
$X\geq\delta I$ para algún $\delta>0$, en una fibra masiva. Los operadores
\begin{equation}
 M=m_{\mathrm{clk}}X,\qquad
 \Lambda=R_{108}X^{-1},\qquad R_g=R_{108}X
 \label{md:hib:operadores-radiales}
\end{equation}
satisfacen $R_g\Lambda=\Lambda R_g=R_{108}^2I$ en todo el espacio.
La afirmación utiliza el cálculo funcional del mismo operador. El núcleo
nulo queda fuera del inverso. El enunciado presente conserva el régimen
acotado uniformemente positivo; la extensión no acotada exige demostrar
por separado los dominios, los cortes compatibles y el paso al límite.
La sección~\ref{md:rigidez:variaciones} demuestra además la conservación
de ambas identidades bajo variaciones no conmutativas y reducción
compatible de memoria.

La interpretación propuesta caracteriza el factor másico relativo como
parámetro de dilatación recíproca de estas dos lecturas. La condición circular
pertenece al resultado. Las longitudes así relacionadas no se identifican
por esta sola igualdad con radios materiales medidos ni con un horizonte
electrónico. La identidad radial aislada se utiliza por su función en la
composición, sin atribuirle prioridad histórica independiente.

\subsection{Respuesta constitutiva, orientación y magnitud}
\label{md:hib:constitutivas}

Con $T=\exp(-xI-yR)$, $R=P_+-P_-$ y proyectores de hoja conservados,
el lector temporal es
\begin{equation}
 \begin{aligned}
 V&=(I-T^{90})(I-T^{120})^{-1}=r_+P_++r_-P_-,\\
 r_\pm&=f(s_\pm),\qquad
 s_\pm=\exp[-30(x\pm y)],\qquad
 f(s)=\frac{1-s^3}{1-s^4}.
 \end{aligned}
 \label{md:hib:operador-constitutivo}
\end{equation}
Se recuperan $\widehat\varepsilon=r_+^2$, $\widehat\mu=r_-^2$,
$\widehat Z=r_-/r_+$ y $\widehat c=(r_+r_-)^{-1}$.
El enlace constitutivo de velocidad establece
$c_{\mathrm{int}}=c_{\mathrm{vac}}=\widehat c\,u_C$ en la misma coordenatización.
La colección orientada de Catalán participa en $f$ mediante su momento
resolvente, conservando la estructura anterior al valor extremo.

La descomposición exacta
\begin{equation}
 V=\widehat c^{-1/2}
 \exp\!\left[-\frac12(\log\widehat Z)R\right]
 \label{md:hib:descomposicion-v}
\end{equation}
distingue dilatación común y reparto entre hojas. El intercambio de hojas
conserva $\widehat c$ e invierte $\widehat Z$; el funcional de Barbero cambia
de signo. La sección~\ref{md:compos:difeomorfismo} demuestra que
$(\widehat c,\gamma)$ recupera $r_+,r_-$ en la colección normalizada.
Su restricción a la imagen generada conserva esta recuperabilidad; el
enunciado de inversión no declara generada toda pareja sintética.

\subsubsection{Convexidad y evaluación de los canales separados}
\label{md:hib:convexidad}

\begin{proposition}
\label{md:hib:proposicion-convexidad}
Para $x>0$ fijo y $|y|<x$, la función
$y\mapsto\log\widehat c(x,y)$ es par y estrictamente convexa.
En consecuencia,
\begin{equation}
 \widehat c(x,y)\geq\widehat c(x,0)
 =\frac{1}{f(e^{-30x})^2},
 \label{md:hib:cota-velocidad}
\end{equation}
con igualdad exactamente para $y=0$.
\end{proposition}
\begin{proof}
Sea $g(u)=\log f(e^{-u})$, para $u>0$. Se tiene
\begin{equation}
 g'(u)=\frac{3}{e^{3u}-1}-\frac{4}{e^{4u}-1}>0,
 \label{md:hib:derivada-g}
\end{equation}
\begin{equation}
 g''(u)=-\frac{9}{4\sinh^2(3u/2)}
 +\frac{16}{4\sinh^2(2u)}<0.
 \label{md:hib:segunda-g}
\end{equation}
La primera desigualdad procede del decrecimiento de
$a/(\exp(au)-1)$ en $a>0$. Para la segunda, la función
$a/\sinh(au/2)$ decrece estrictamente: su derivada tiene el signo de
$\sinh v-v\cosh v$, negativo para $v>0$. Como
\[
 \log\widehat c(x,y)=-g(30(x+y))-g(30(x-y)),
\]
su segunda derivada respecto de $y$ es
\begin{equation}
 -900\bigl[g''(30(x+y))+g''(30(x-y))\bigr]>0.
 \label{md:hib:segunda-velocidad}
\end{equation}
La paridad fija el mínimo único en cero.
\end{proof}
También se obtiene
\begin{equation}
 \begin{aligned}
 \log\widehat c(x,y)
 &=-2g(30x)-900g''(30x)y^2+O(y^4),\\
 \log\widehat Z(x,y)
 &=-60g'(30x)y+O(y^3).
 \end{aligned}
 \label{md:hib:desarrollos-canales}
\end{equation}
La velocidad pierde el signo de orientación, pero conserva una respuesta
cuadrática a la separación de canales. La impedancia conserva información
orientada desde el primer orden. Promediar $x+y$ y $x-y$ antes de aplicar el
lector elimina aquella contribución. La afirmación es una consecuencia
demostrada de la respuesta constitutiva, distinta de una medición nueva.

Sobre la rama positiva de acción,
$\eta_{\mathrm{ret}}=(90/\pi)(y-x/500)$; por tanto,
\begin{equation}
 \hbar_{\mathrm{ret}}=s_0\frac{90}{\pi}(y-x/500),\qquad
 E_k=s_0\frac{90u_C|k|}{\pi}(y-x/500)\widehat c(x,y).
 \label{md:hib:energia-angular}
\end{equation}
Para $x$ fijo, $|k|>0$ y $x/500<y<x$, esta última función crece
estrictamente en $y$. Es una propiedad de la colección de realización
radiativa homogénea. Su aplicación a una trayectoria generada conserva el
constructor de $y$; la variación de un parámetro de esta colección no equivale
por sí sola a una variación experimental de una constante universal.
Bajo inversión de orientación se transporta también la etiqueta de acción,
manteniendo la rama positiva correspondiente.

\subsection{Acción, reloj e información térmica}
\label{md:hib:termica}

La sección térmica conserva
$k_B\Theta_{108}\log3=E_{\mathrm{clk}}$, con
$\Theta_{108}=\Theta_{\mathrm{clk,ret}}$ en la sección de retorno.
Definimos
$\beta_{108}=(k_B\Theta_{108})^{-1}$. Para
$E_\Gamma=E_{\mathrm{clk}}\Xi_\Gamma$ se deduce exactamente
\begin{equation}
 \beta_{108}E_\Gamma=(\log3)\Xi_\Gamma,
 \qquad
 \boxed{\exp(-\beta_{108}E_\Gamma)=3^{-\Xi_\Gamma}.}
 \label{md:hib:peso-termico}
\end{equation}
La expresión es un peso de Gibbs no normalizado. Para una colección finita
de estados, o una colección con suma de partición convergente, y
multiplicidades $g_\Gamma$, se obtiene
\begin{equation}
 Z_{\mathrm{part}}=\sum_\Gamma g_\Gamma3^{-\Xi_\Gamma},\qquad
 p_\Gamma=\frac{g_\Gamma3^{-\Xi_\Gamma}}{Z_{\mathrm{part}}},\qquad
 \frac{p_Y/g_Y}{p_X/g_X}=3^{-(\Xi_Y-\Xi_X)}.
 \label{md:hib:particion}
\end{equation}
Se asume una realización canónica a $\Theta_{108}$ y un origen energético
común. Cuando interviene un potencial químico, se incorpora a la energía
correspondiente. El índice $\Gamma$ enumera estados o grupos físicos
distintos; varias historias equivalentes no producen automáticamente una
multiplicidad. Las entropías adimensionales son
\begin{equation}
 \begin{aligned}
 H_{\mathrm{micro}}
 &=\log Z_{\mathrm{part}}+(\log3)\sum_\Gamma p_\Gamma\Xi_\Gamma,\\
 H_{\mathrm{grupos}}
 &=H_{\mathrm{micro}}-\sum_\Gamma p_\Gamma\log g_\Gamma.
 \end{aligned}
 \label{md:hib:entropias}
\end{equation}
Para afirmar la finitud de la entropía se requiere, además,
\begin{equation}
 \sum_\Gamma g_\Gamma\Xi_\Gamma3^{-\Xi_\Gamma}<\infty.
 \label{md:hib:condicion-entropia}
\end{equation}
La finitud de $Z_{\mathrm{part}}$ garantiza por sí sola la normalización.
En colecciones finitas se cumplen ambas condiciones.

La fórmula enlaza la estructura másica con frecuencia propia y ponderación
térmica en el mismo marco. Conserva la distinción entre energía y calor,
entre coordenada real $\Xi_\Gamma$ y conteo entero de trits, y entre una
realización a $\Theta_{108}$ y una afirmación sobre la temperatura del
universo. Para ocupaciones bosónicas o fermiónicas libres por modo, con
potencial químico nulo, la misma sustitución produce
\begin{equation}
 \overline n_\Gamma=\frac{1}{3^{\Xi_\Gamma}\mp1}.
 \label{md:hib:ocupaciones}
\end{equation}
El signo menos corresponde a Bose y el signo más a Fermi. El total del
grupo degenerado es $g_\Gamma\overline n_\Gamma$. En Bose se exige
$\Xi_\Gamma>0$ y el modo nulo se trata separadamente. Son realizaciones
estadísticas posteriores aplicadas a un espectro ya construido.

\subsubsection{Compatibilidad conjunta de cuatro lecturas}
\label{md:hib:cuatro-lecturas}

Componiendo las secciones anteriores, bajo sus mismas hipótesis,
\begin{equation}
 \boxed{\Xi_\Gamma=\frac{\Omega_\Gamma}{\omega_0}
 =\frac{r_{g,\Gamma}}{R_{108}}
 =\frac{R_{108}}{\bar\lambda_\Gamma}
 =\frac{E_\Gamma}{k_B\Theta_{108}\log3}.}
 \label{md:hib:compatibilidad-cuatro}
\end{equation}
En la colección circular, $\Xi=1$ identifica la sección de frecuencia
$\omega_0$, radios coincidentes $R_{108}$ y energía $E_{\mathrm{clk}}$.
Su peso térmico no normalizado es $1/3$. Es un punto de compatibilidad de
las lecturas. La existencia de una especie del atlas con $\Xi=1$, o de
una segunda especie generada para cada sección recíproca $1/\Xi$, requiere
su pertenencia producida al atlas; no se infiere de esta igualdad.

\subsection{Identidad energética de la elipse y transporte exacto}
\label{md:hib:elipse}

En la elipse de acción,
\begin{equation}
 a_S=\sqrt{2\hbar}\,e^{\zeta/2},\qquad
 b_S=\sqrt{2\hbar}\,e^{-\zeta/2},\qquad
 Q=a_S\cos\theta,\qquad P=b_S\sin\theta,
 \label{md:hib:elipse-coordenadas}
\end{equation}
con $\zeta=\operatorname{artanh}(C^*_{\deg}/A_{\deg})$, se tiene
$a_Sb_S=2\hbar$, área orientada
$\mathcal A(\theta)=\hbar\theta$ y $\mathcal A(2\pi)=h$.
La fase paramétrica $\theta$ y el ángulo polar de un punto de la elipse
conservan sus significados geométricos respectivos.

La realización LC declarada utiliza
\begin{equation}
 L_\zeta=\frac{Z_*}{\omega}e^{-\zeta},\qquad
 C_\zeta=\frac{1}{Z_*\omega}e^\zeta.
 \label{md:hib:lc}
\end{equation}
En coordenadas de acción $Q=\sqrt{Z_*}\,q$ y
$P=\Phi/\sqrt{Z_*}$,
\begin{equation}
 H_0=\frac{\omega}{2}
 \left(e^{-\zeta}Q^2+e^\zeta P^2\right),\qquad
 U_E=\hbar\omega\cos^2\theta,\qquad
 U_B=\hbar\omega\sin^2\theta.
 \label{md:hib:energia-elipse}
\end{equation}
Se obtiene $H_0=\hbar\omega$ por composición de las dos contribuciones.
El reloj y el transductor pertenecen a la realización declarada; la sección
$\hbar$ conserva su construcción anterior al circuito. La impedancia
$Z_{\mathrm{LC}}/Z_*=e^{-\zeta}$ y la impedancia constitutiva del vacío
$f(s_-)/f(s_+)$ son lectores distintos. Su antecedente angular común no
establece una igualdad entre ambos.

\subsubsection{Transporte discreto de la forma cuadrática}
\label{md:hib:transporte-discreto}

\begin{proposition}
\label{md:hib:proposicion-discreta}
Sean $S_\zeta=\operatorname{diag}(e^{\zeta/2},e^{-\zeta/2})$,
$g_\zeta=S_\zeta^{-T}S_\zeta^{-1}$ y
$\mathsf R(-\theta)$ una rotación. Para valores de forma $\zeta_n$ y
fases $\theta_n$, el transporte
\begin{equation}
 X_{n+1}=S_{\zeta_{n+1}}\mathsf R(-\theta_n)S_{\zeta_n}^{-1}X_n
 =T_nX_n
 \label{md:hib:transporte}
\end{equation}
conserva exactamente
$I_n=X_n^Tg_{\zeta_n}X_n/2$ y tiene determinante uno.
\end{proposition}
\begin{proof}
La sustitución y la identidad $\mathsf R^T\mathsf R=I$ dan
\begin{equation}
 T_n^Tg_{\zeta_{n+1}}T_n=g_{\zeta_n}.
 \label{md:hib:isometria-transportada}
\end{equation}
La forma cuadrática se conserva al evaluar ambos lados sobre $X_n$.
Los factores de $T_n$ tienen determinante uno, de donde
$\det T_n=1$.
\end{proof}
Es una composición del transporte radial previamente construido. Si
$\zeta_n,\theta_n$ se extraen de una ruta TPK, la identidad se aplica a
esa secuencia; la proposición algebraica conserva separada la selección
de la ruta.

\subsubsection{Transporte diferenciable y término de conexión}
\label{md:hib:conexion}

Si $\zeta(t)$ es de clase $C^1$ y $\omega(t)>0$ es continua, la relación
$X=S_\zeta Y$ y la evolución $\dot Y=-\omega J_0Y$, con
\begin{equation}
 J_0=\begin{pmatrix}0&-1\\1&0\end{pmatrix},
 \label{md:hib:j-canonico}
\end{equation}
dan
\begin{equation}
 \dot Q=\frac{\dot\zeta}{2}Q+\omega e^\zeta P,\qquad
 \dot P=-\omega e^{-\zeta}Q-\frac{\dot\zeta}{2}P.
 \label{md:hib:ecuaciones-transporte}
\end{equation}
Con las convenciones $\dot Q=\partial_PH$ y
$\dot P=-\partial_QH$, el Hamiltoniano es, salvo un sumando dependiente
sólo del tiempo,
\begin{equation}
 \boxed{H_{\mathrm{tot}}=
 \frac{\omega}{2}\left(e^{-\zeta}Q^2+e^\zeta P^2\right)
 +\frac{\dot\zeta}{2}QP.}
 \label{md:hib:hamiltoniano-conexion}
\end{equation}
El invariante
\begin{equation}
 I=\frac12\left(e^{-\zeta}Q^2+e^\zeta P^2\right)
 \label{md:hib:invariante}
\end{equation}
se conserva exactamente, sin requerir una aproximación adiabática.
\begin{proof}
Al derivar se obtiene
\[
 \dot I=\frac{\dot\zeta}{2}
 \left(e^\zeta P^2-e^{-\zeta}Q^2\right)
 +e^{-\zeta}Q\dot Q+e^\zeta P\dot P.
\]
La sustitución de \eqref{md:hib:ecuaciones-transporte} cancela los términos
proporcionales a $\dot\zeta$ y los dos términos $\omega QP$.
Así, $\dot I=0$. Si se omite el término mixto de
\eqref{md:hib:hamiltoniano-conexion}, queda
\begin{equation}
 \dot I=\frac{\dot\zeta}{2}
 \left(e^\zeta P^2-e^{-\zeta}Q^2\right).
 \label{md:hib:defecto-sin-conexion}
\end{equation}
El signo mixto contrario duplica ese defecto.
\end{proof}
Sobre la elipse, $QP=I\sin(2\theta)$ y $\dot\theta=-\omega$, por lo que
\begin{equation}
 H_{\mathrm{tot}}=I\omega+\frac{I\dot\zeta}{2}\sin(2\theta).
 \label{md:hib:hamiltoniano-elipse}
\end{equation}
La matriz de la forma cuadrática tiene determinante
$\omega^2-\dot\zeta^2/4$. Para $\omega>0$, es definida positiva
exactamente cuando $|\dot\zeta|<2\omega$, y semidefinida positiva en
igualdad. Este umbral concierne a la positividad: la existencia del
transporte y la estabilidad física conservan sus condiciones respectivas.
La conservación de $I$ no implica conservación de $H_{\mathrm{tot}}$,
que depende explícitamente del tiempo.

La versión cuántica exige la simetrización del término mixto,
$\dot\zeta(QP+PQ)/4$. Esta expresión indica la operación requerida sin
presuponer una representación cuántica adicional.

\subsubsection{Compatibilidad con la sección de acción generadora del contraángulo}
\label{md:hib:compatibilidad-accion}

En la construcción rectora,
\begin{equation}
 \zeta=\operatorname{artanh}
 \left[\frac{2(\eta_{\mathrm{ret}}+\alpha)}{1000\alpha}\right].
 \label{md:hib:forma-genealogica}
\end{equation}
Fijar $\alpha$ y $\eta_{\mathrm{ret}}$ fija también $\zeta$.
El resultado con $\zeta(t)$ variable es, por tanto, una extensión de
transporte entre formas. Para representar una trayectoria física del par
rector, la variación debe proceder de la genealogía correspondiente, y no
de una deformación elegida de manera independiente.

También puede componerse exactamente el transporte cuando cambia la
sección positiva de acción entre dos estados, conservando la misma coordenatización
dimensional. Sean
\begin{equation}
 \hbar_n>0,\qquad
 \rho_{\mathrm{act},n}=\frac{\hbar_{n+1}}{\hbar_n},\qquad
 \widetilde T_n=\sqrt{\rho_{\mathrm{act},n}}\,
 S_{\zeta_{n+1}}\mathsf R(-\theta_n)S_{\zeta_n}^{-1}.
 \label{md:hib:transporte-conforme}
\end{equation}
Entonces
\begin{equation}
 \det\widetilde T_n=\rho_{\mathrm{act},n},\qquad
 \widetilde T_n^Tg_{\zeta_{n+1}}\widetilde T_n
 =\rho_{\mathrm{act},n}g_{\zeta_n},\qquad
 \frac{I_{n+1}}{\hbar_{n+1}}=\frac{I_n}{\hbar_n}.
 \label{md:hib:accion-normalizada}
\end{equation}
Las igualdades siguen de \eqref{md:hib:isometria-transportada} al multiplicar
por el factor escalar. Es un transporte conformemente simpléctico: conserva
la acción normalizada y conserva el área de acción sin normalizar
exactamente cuando $\rho_{\mathrm{act},n}=1$.
Su límite diferenciable tiene traza $\dot\hbar/\hbar$. Un generador
hamiltoniano lineal en un plano con forma simpléctica fija tiene traza cero;
por ello, una variación material de $\hbar$ exige especificar la forma
transportada o los grados de libertad con los que se intercambia acción.
El término mixto de deformación, por sí solo, conserva traza cero.

El resultado es un criterio de compatibilidad de las realizaciones. Para
interpretar variaciones numéricas debidas a un cambio de unidades se
transportan primero las bases: el cambio de coordenatización mantiene la sección
física. La composición no afirma una variabilidad experimental de la
constante de Planck.

\subsubsection{Transformación de representación y deformación material}
\label{md:hib:interpretacion-transporte}

Si $S_\zeta$ es un cambio pasivo de coordenadas, el término mixto es una
conexión de representación. Se transportan también los observables,
\begin{equation}
 L_X(S_\zeta Y)=L_Y(Y).
 \label{md:hib:observables-transportados}
\end{equation}
Esta igualdad conserva la lectura física al cambiar la expresión del
Hamiltoniano, sin introducir una interacción adicional ni trabajo físico.

Una deformación material exige obtener $\zeta$ de la actualización TPK y
mostrar un cambio relativo entre sistema y referencia, ambos transportados
coherentemente. Puede contrastarse una fase relativa, una respuesta
$Z_{\mathrm{sistema}}/Z_{\mathrm{referencia}}$ o una transferencia entre
modos definidos por el detector. Deben distinguirse dos historias TPK no
equivalentes por cambio de representación y conservar sus condiciones de
lectura. Cuando se atribuya un intercambio energético, el balance comprende
campo, memoria y dispositivo. El teorema de transporte delimita esta
distinción y conserva separada la ley dinámica que la haga física.

\subsection{Momento, contrastes y alcance de la composición}
\label{md:hib:momento}

La realización de Clifford--Dirac recibe la masa de ruta. En una misma
sección de acción, su coeficiente se escribe
\begin{equation}
 \frac{m_\Gamma c}{\hbar}
 =\frac{\omega_0\Xi_\Gamma}{c}.
 \label{md:hib:coeficiente-dirac}
\end{equation}
En una realización libre y homogénea, con la relación relativista de
dispersión declarada, ésta se reexpresa como
\begin{equation}
 \Omega^2=c^2|k|^2+\omega_0^2\Xi_\Gamma^2.
 \label{md:hib:dispersion}
\end{equation}
La reescritura conserva las hipótesis de esa realización y sus antecedentes
de conexión espacio-temporal. Los momentos de retorno $q_-^h-q_+^h$
son evaluaciones de las torres y se distinguen de los momentos mecánicos.

El factor $\Xi_\Gamma$ reaparece como peso másico, frecuencia relativa,
dilatación radial y exponente térmico. La respuesta constitutiva determina
la conversión espacial común. Barbero conserva información orientada
complementaria a la velocidad. La acción fija la medida areal; el reloj
introduce la tasa temporal; la forma y su transporte especifican el reparto
y su evolución.

Los contrastes correspondientes comprenden cocientes de masas y frecuencias;
compatibilidad de velocidad, impedancia y Barbero; conservación de la forma
cuadrática bajo transporte; y pesos térmicos una vez fijadas temperatura y
multiplicidades. Estas relaciones preservan más datos que una coincidencia
decimal aislada. Sus pruebas mantienen el balance energético y las
condiciones de las entropías consideradas; no implican entropía
termodinámica nula.
